% ============================================================================
% R\'ESUM\'E FRAN\c{C}AIS
% Fichier : preambule/resume.tex
% ============================================================================

\chapter*{R\'ESUM\'E}
\markboth{R\'esum\'e}{}
\addcontentsline{toc}{chapter}{R\'esum\'e}
\thispagestyle{preliminaire}
\begin{sloppypar}
La cryptographie classique appr\'ecie la s\'ecurit\'e d'une preuve \`a
l'aune de sa robustesse face \`a des adversaires computationnellement
born\'es. \`A l'inverse, le paradigme institutionnel impose que la preuve
demeure intelligible pour un arbitre non sp\'ecialiste afin d'\'eclairer
une d\'ecision. Les d\'eploiements op\'erationnels r\'ev\`elent une
pathologie r\'ecurrente. Une preuve irr\'eprochable du point de vue
math\'ematique peut se r\'ev\'eler inutilisable par l'agent charg\'e de
trancher. Cette tension structurelle entre les deux paradigmes constitue
le paradoxe de l'authenticit\'e invisible, selon lequel plus une preuve
prot\`ege efficacement, plus elle devient opaque pour un tiers profane.

La th\`ese introduit une m\'etrique informationnelle in\'edite,
l'opposabilit\'e Opp, d\'efinie comme la proportion d'incertitude
institutionnelle qu'une preuve parvient \`a dissiper. Elle propose le
cadre CLO, qui agr\`ege confidentialit\'e (Conf), fiabilit\'e (Rel) et
opposabilit\'e (Opp) et l'articule \`a des param\`etres cognitifs
explicites, en particulier la capacit\'e d'extraction institutionnelle
$\Delta$ et la granularit\'e des verdicts $\alpha$. Elle \'etablit un
trilemme d'impossibilit\'e conditionnel. Sous quatre hypoth\`eses bornant
les capacit\'es de l'institution, le produit Conf$\cdot$Rel$\cdot$Opp
est major\'e par $\Delta/(\gamma\lambda)$, o\`u $\gamma = \alpha H(C) - 1$,
\`a un terme d'ordre $\lambda^{-2}$ et \`a un reste n\'egligeable pr\`es.
Une forme raffin\'ee identifie le r\'egime o\`u la capacit\'e
institutionnelle devient la contrainte liante. La ressource limitante
est alors la quantit\'e d'information qu'une institution peut
extraire humainement d'une preuve, bien plus que la taille des cl\'es.

Ce r\'esultat se prolonge au contexte post-quantique. Dans le sc\'enario
d'un support ou d'un v\'erificateur quantique, le bruit induit une perte
d'interpr\'etabilit\'e $\eta_q$, nulle pour les preuves classiques, dont
l'effet sur la borne cro\^it avec l'entropie contextuelle. La th\`ese
propose aussi la famille SOTP, param\'etr\'ee pour diff\'erents
r\'egimes de richesse contextuelle, dont elle \'etablit les propri\'et\'es
cryptographiques. Son optimalit\'e asymptotique est formul\'ee comme
conjecture. Le noyau alg\'ebrique du
trilemme est certifi\'e m\'ecaniquement dans Lean4, qui explicite
hypoth\`eses, constantes et domaine de validit\'e et fournit deux
t\'emoins formels garantissant la coh\'erence interne et la port\'ee du
r\'esultat.

Enfin, la th\`ese sp\'ecifie un programme exp\'erimental r\'efutable
comprenant six protocoles destin\'es \`a mesurer les param\`etres
institutionnels, avec estimateurs et seuils de r\'efutation, puis une
calibration illustrative de ces param\`etres sur des syst\`emes
d\'eploy\'es. Elle v\'erifie la plausibilit\'e des ordres de grandeur sans
constituer une validation empirique, le programme n'ayant pas \'et\'e ex\'ecut\'e en conditions
prospectives.

En \'erigeant l'opposabilit\'e au rang de propri\'et\'e cryptographique de
premi\`ere classe, ce travail constitue la cryptographie institutionnelle
comme domaine th\'eorique autonome. Il fournit des instruments d'aide au
choix des protocoles dans des secteurs sensibles tels que le vote
\'electronique, la conformit\'e r\'eglementaire et la transition
post-quantique et pr\'epare une normalisation future int\'egrant
l'interpr\'etabilit\'e parmi les crit\`eres de s\'ecurit\'e.

\vspace{0.5cm}
\noindent\textbf{Mots-cl\'es~:} cryptographie institutionnelle,
opposabilit\'e, Trilemme CRO, protocoles SOTP, v\'erification formelle,
programme exp\'erimental falsifiable, Lean4.
\end{sloppypar}